\chapter{Internal Block Diagram}

\begin{lernziele}
Nach diesem Kapitel können Sie Datenflüsse und interne Verbindungen zwischen Komponenten im IBD darstellen.
\end{lernziele}

\section{Zweck des IBD}
Das \gls{ibd} zeigt, wie Teile eines Blocks intern verbunden sind. Während das BDD sagt, welche Bausteine existieren, zeigt das IBD, wie sie zusammenarbeiten.

\section{Beispiel Datenfluss}
\begin{verbatim}
Camera -> ObjectDetector -> NavigationSystem
Lidar  -> SLAM           -> NavigationSystem
NavigationSystem -> MotionControl -> Motors
BatterySystem -> NavigationSystem
\end{verbatim}

\section{Ports}
Ports beschreiben Ein- und Ausgänge von Blöcken. Ein Kamerablock kann einen Port \texttt{imageOut} haben. Der ObjectDetector kann einen Port \texttt{imageIn} und \texttt{objectsOut} haben.

\section{Schnittstellen}
Eine Verbindung ohne klare Schnittstelle ist wenig wert. Deshalb sollte dokumentiert werden, welche Daten übertragen werden. Beispiel: \texttt{sensor\_msgs/Image} für Kamerabilder in ROS2.

\section{IBD und ROS2}
Das IBD passt gut zur ROS2-Kommunikation. ROS2-Topics können als Datenflüsse modelliert werden. Ein Topic \texttt{/camera/image\_raw} verbindet \texttt{camera\_node} und \texttt{object\_detector\_node}.

\begin{uebungbox}
Erstellen Sie ein IBD für das \texttt{PerceptionSystem}. Modellieren Sie Kamera, Lidar, ObjectDetector und ObstacleDetector inklusive Datenflüssen.
\end{uebungbox}
